\documentclass[10pt]{article}
\usepackage[T1]{fontenc}
\usepackage{lmodern,graphicx,amsmath,booktabs,longtable}
\usepackage[a4paper,margin=24mm]{geometry}
\usepackage[hidelinks]{hyperref}
\graphicspath{{build/}}
% GENERATED by reproducibility/make_numbers.py -- do not edit.
% Sources:
%   bench/results/wp8-batch-graniterapids.json
%   bench/results/wp8-batch-tremont-20260820t105514z.json
%   bench/results/realdock-1ggr-f32-tremont-r5-nvidia-geforce-rtx-4090-gpu.json
%   bench/results/realdock-1ggr-f64-tremont-r5-nvidia-geforce-rtx-4090-gpu.json
%   results/wp8/comparison_summary.csv

% ---- the benchmark campaign, first run ----
\newcommand{\campCases}{52}
\newcommand{\campOk}{51}
\newcommand{\campRankOne}{27}
\newcommand{\campRankOnePct}{52.9}
\newcommand{\campHitTopK}{44}
\newcommand{\campHitTopKPct}{86.3}
\newcommand{\campClusterOne}{15}
\newcommand{\campClusterOnePct}{29.4}
\newcommand{\campIrmsdFour}{27}
\newcommand{\campSeconds}{1909.0}
\newcommand{\campMinutes}{31.8}
\newcommand{\campPerCaseLo}{34.8}
\newcommand{\campPerCaseHi}{45.1}
\newcommand{\campFloorLo}{0.42}
\newcommand{\campFloorHi}{1.20}
\newcommand{\campFloorMed}{0.73}
\newcommand{\campAngleFloor}{2.37}
\newcommand{\campNrot}{69{,}960}
\newcommand{\campWone}{7}
\newcommand{\campHitAll}{50}
\newcommand{\campMaxHitRank}{24{,}544}
\newcommand{\campEqualRanks}{51}
\newcommand{\campEqualEnergy}{35}
\newcommand{\dockBestRank}{293}
\newcommand{\campTableBody}{1ACB & 1 & 206 & 1 & 40.9 \\
1AVX & 19 & 42 & 7 & 40.8 \\
1AY7 & 490 & 2 & 26 & 34.9 \\
1BVN & 1 & 351 & 1 & 36.9 \\
1CGI & 1 & 567 & 1 & 38.6 \\
1CLV & 1 & 465 & 1 & 35.2 \\
1D6R & -- & 0 & -- & 36.0 \\
1DFJ & 1 & 16 & 17 & 35.5 \\
1E6E & 3 & 139 & 2 & 34.8 \\
1EAW & 95 & 9 & 11 & 37.3 \\
1EWY & 204 & 23 & 8 & 35.8 \\
1EZU & 3 & 67 & 3 & 39.1 \\
1F34 & 1 & 130 & 2 & 36.0 \\
1F6M & 9 & 64 & 5 & 39.6 \\
1FLE & 1 & 255 & 1 & 36.9 \\
1FQ1 & 4938 & 0 & -- & 36.5 \\
1GL1 & 1 & 175 & 2 & 35.6 \\
1GXD & 25 & 31 & 9 & 36.3 \\
1HIA & 1 & 186 & 1 & 36.4 \\
1IJK & 121 & 16 & 12 & 35.8 \\
1JIW & 22 & 8 & 21 & 37.8 \\
1JTG & 1 & 76 & 2 & 36.3 \\
1KKL & 1615 & 0 & -- & 36.5 \\
1M10 & 451 & 1 & -- & 38.6 \\
1MAH & 1 & 78 & 6 & 41.3 \\
1N8O & 1 & 150 & 2 & 38.6 \\
1NW9 & 1 & 215 & 1 & 41.0 \\
1OC0 & 2 & 187 & 11 & 40.6 \\
1OPH & 19 & 45 & 6 & 40.3 \\
1OYV\_A & 1 & 134 & 2 & 36.5 \\
1OYV\_B & 5130 & 0 & -- & 45.1 \\
1PPE & 1 & 172 & 3 & 38.5 \\
1PXV & 1 & 31 & 11 & 37.0 \\
1R0R & 1 & 147 & 4 & 36.2 \\
1TMQ & 1 & 209 & 2 & 35.5 \\
1UDI & 1 & 245 & 1 & 38.6 \\
1YVB & 1 & 151 & 1 & 36.4 \\
1ZLI & 1 & 174 & 1 & 36.4 \\
2ABZ & 19 & 105 & 4 & 39.1 \\
2B42 & 1 & 39 & 7 & 35.7 \\
2J0T & 1 & 152 & 1 & 38.6 \\
2MTA & 4 & 254 & 1 & 35.6 \\
2O3B & 2948 & 0 & -- & 35.8 \\
2O8V & 1 & 196 & 2 & 35.6 \\
2OUL & 1 & 163 & 1 & 36.2 \\
2PCC & 5720 & 0 & -- & 39.0 \\
2SIC & 1 & 279 & 1 & 38.6 \\
2SNI & 1 & 211 & 1 & 35.9 \\
2UUY & 927 & 1 & -- & 37.5 \\
3SGQ & 376 & 8 & 10 & 35.5 \\
4CPA & \multicolumn{4}{l}{Atom typing failed} \\
7CEI & 24544 & 0 & -- & 35.8 \\}

% ---- the second campaign, a different machine ----
\newcommand{\campTwoSeconds}{2563.5}
\newcommand{\campTwoMinutes}{42.7}
\newcommand{\campTwoRankOne}{27}
\newcommand{\campTwoHitTopK}{44}
\newcommand{\campTwoClusterOne}{15}

% ---- the real complex, both arms in one process ----
\newcommand{\dockGpuFThirtyTwo}{20.7}
\newcommand{\dockGpuFSixtyFour}{60.5}
\newcommand{\dockCpuFThirtyTwo}{1083.7}
\newcommand{\dockSpeedupFThirtyTwo}{52.3}
\newcommand{\dockSpeedupFSixtyFour}{18.7}
\newcommand{\dockAtoms}{1{,}774}
\newcommand{\dockNrot}{69{,}960}
\newcommand{\dockRank}{1}
\newcommand{\dockRmsd}{8.911}
\newcommand{\dockBest}{0.85974}
\newcommand{\dockFloor}{0.85974}
\newcommand{\dockEnergy}{-659.041}
\newcommand{\parityCompared}{2{,}000}
\newcommand{\parityCellAgree}{1.0}
\newcommand{\parityRelEnergy}{5.9e-06}

% ---- the external comparison, five cases ----
\newcommand{\wpCases}{5}
\newcommand{\wpJoint}{0}
\newcommand{\wpSameBand}{3}
\newcommand{\wpNearN}{4}
\newcommand{\wpNearLo}{1.85}
\newcommand{\wpNearHi}{1.98}
\newcommand{\wpNearSpread}{0.13}
\newcommand{\wpRankLo}{20}
\newcommand{\wpRankHi}{437}
\newcommand{\wpCentreLo}{3.77}
\newcommand{\wpCentreHi}{6.20}
\newcommand{\wpCentreAllLo}{3.77}
\newcommand{\wpCentreAllHi}{6.68}
\newcommand{\wpOddCase}{1AY7}
\newcommand{\wpOddNear}{5.76}
\newcommand{\wpTableBody}{1ACB & 1 & 1 & 4.58 & 1 & 1.93 & 319 \\ 1AVX & 1 & 7 & 5.60 & 7 & 1.85 & 20 \\ 1AY7 & -- & 26 & 6.68 & 7 & 5.76 & 338 \\ 1BVN & 1 & 1 & 3.77 & 1 & 1.98 & 106 \\ 1CGI & 1 & 1 & 6.20 & 1 & 1.98 & 437 \\}

% ---- carried by hand: no committed JSON holds these ----
%   tests on CPU                       5,205                          <- console output of Pkg.test()
%   tests with the GPU extension       5,282                          <- console output of a device deploy
%   correlation vs direct O(N^6)       5.1e-16                        <- test_correctness.jl
%   CPU/GPU agreement, Float32         3e-7                           <- device suite, paper1_hpc 6.5.1
%   CPU/GPU agreement, Float64         7e-16                          <- device suite, paper1_hpc 6.5.1
%   pipeline vs its own floor, CPU     1.11-1.52x                     <- paper1_hpc 6.3
%   pipeline vs its own floor, GPU     1.17-2.50x                     <- paper1_hpc 6.5.2
%   Fourier-domain accumulation        1.56-1.97x                     <- paper1_hpc 6.4 row A1
%   real-input transforms              1.44-3.29x                     <- paper1_hpc 6.4 row A2
%   translation scan rewrite           16.5-60.3x                     <- paper1_hpc 6.5.5
%   closed-form pose RMSD              1225x at M=2000                <- paper1_hpc 2.6
%   shape-weight crossover             3.264 / 3.151                  <- paper1_hpc 6.8
%   fibre pruning                      0 of 200, 1.34x slower         <- paper1_hpc 7.1
%   restraints                         78.06% infeasible, 1.26-1.86x  <- paper1_hpc 9.1
%   Cn symmetry                        69.4x                          <- paper1_hpc 9.1
%   noise floor                        9.5-21.5% / 1.4% / 1.6-4.7%    <- paper1_hpc 6.1

\renewcommand{\thetable}{S\arabic{table}}
\renewcommand{\thefigure}{S\arabic{figure}}
\begin{document}
\begin{center}
{\LARGE Supplementary Information for JuPIPER\par}
\medskip
{\large Editable scoring for FFT-based protein--protein docking on CPUs and GPUs\par}
\medskip
Ilya Afanasyev, Innopolis University
\end{center}
\section{Correlation on a translation grid}
The receptor and ligand are represented on a Cartesian grid. Let $R_p$ and
$L_p$ be the fields for score component $p$, including its weight on the
receptor side. The discrete correlation gives the energy at every translation:
\[
E(\mathbf{t})=\sum_p\sum_{\mathbf{x}}R_p(\mathbf{x})L_p(\mathbf{x}+\mathbf{t})
=\mathcal F^{-1}\!\left(\sum_p\overline{\mathcal F(R_p)}\mathcal F(L_p)\right)(\mathbf{t}).
\]
For $m$ fields, each orientation requires $m$ forward transforms and one inverse
transform. Receptor spectra are reusable across orientations. Linearity of the
inverse transform allows the products for all fields to be summed in Fourier
space. The index of the minimum of the resulting grid selects a translation.

\begin{figure}[!ht]
\renewcommand{\thefigure}{S\arabic{figure}(b)}
\centering\includegraphics[width=.91\linewidth]{figS1_dispatch.pdf}
\caption{CPU and GPU implementations of the correlation. The numerical
tests compare full synthetic energy fields and the translations selected on 1GGR.}
\end{figure}

\begin{figure}[!ht]
\centering\includegraphics[width=.60\linewidth]{fig3_stages_compact.pdf}
\caption{Synthetic per-orientation stage timings, one virtual CPU core and
one RTX~4090, Float32, $P=6$, batch size one. Each of the five separately
timed stages is divided by their sum. The translation scan follows the
correlation. Forward FFTs lead on the CPU at all three sizes. On the GPU,
voxelisation leads at $64^3$ (40.8\%) and forward transforms lead at $96^3$
and $128^3$ (38.0\% and 40.2\%).}
\end{figure}

\section{Branch-and-bound test on Hopf fibres}
A branch-and-bound experiment grouped orientations by Hopf fibre. Its lower
bound retired zero of 200 fibres at the tested thresholds and increased the
runtime by a factor of 1.34.

\clearpage
\section{Bound redocking by case}
The benchmark category contains \campCases{} cases. Two use different chains
of 1OYV. The control 1GGR is outside the category. A hit is a pose within
10\,\AA{} ligand C$\alpha$ RMSD of the deposited placement. A dash denotes
no qualifying pose or cluster. First-hit rank covers the full search, while
the hit count uses the first 1,000 poses. Times cover the orientation loop
on the first machine. Of the \campOk{} completed searches, \campHitAll{} had
a hit somewhere in the ranking. For 1D6R the stored first-hit rank is zero.

\setcounter{table}{1} % Preserve the benchmark table's identifier, S2.
\small
\begin{longtable}{lrrrr}
\caption{First campaign, including the failed case.}\\
\toprule Case & First-hit rank & Hits in top 1,000 & First near-native cluster & Search (s)\\\midrule\endfirsthead
\toprule Case & First-hit rank & Hits in top 1,000 & First near-native cluster & Search (s)\\\midrule\endhead
\campTableBody
\bottomrule
\end{longtable}
\end{document}
